\section{Introduction}\label{section::introduction}

\todo{Why are LLMs so important right now? which industries are they getting into? how pervasive are they in our lives right now? Have they replaced google?}

Large Language Models (LLMs) have rapidly emerged as transformative technologies, exerting a significant influence across a multitude of industries. 
Their applications are increasingly noted in sectors such as finance \cite{dowling2023chatgpt}, healthcare \cite{thirunavukarasu2023large}, legal services \cite{bommarito2022gpt}, and education \cite{kasneci2023chatgpt, 10822885}. 
The pervasive nature of LLMs is also shown by their growing presence in public discourse and their diverse utility. 
They are employed for tasks ranging from coding assistance \cite{shu2024llms} and as pedagogical tools to support problem-solving in higher education \cite{10.1145/3657604.3662036}, to addressing complex inquiries that necessitate more contextual understanding than conventional search engine queries typically provide \cite{liu2024llm}.
The genesis of modern LLMs can be traced back to advancements within the field of Natural Language Processing (NLP). A pivotal moment was the introduction of the Transformer architecture by Google researchers, initially aimed at enhancing machine translation capabilities \cite{vaswani2017attention}. This architecture has since become the foundational framework for the majority of contemporary LLMs. Subsequent developments, such as Bidirectional Encoder Representations from Transformers (BERT), marked a significant leap, enabling models to be pre-trained on vast amounts of text data and then fine-tuned for a variety of specific downstream tasks \cite{devlin2018bert}. This contrasted with earlier models and paved the way for more versatile and powerful language understanding.

The transition of LLMs from specialized research niches to widespread public awareness was notably accelerated in late 2022 with the release of OpenAI's ChatGPT. Based on the GPT-3.5 architecture (an evolution of the model described in \cite{brown2020languagemodelsfewshotlearners} which focused on GPT-3), ChatGPT garnered immense public interest, reportedly reaching approximately 100 million monthly active users by January 2023 \cite{hu_chatgpt_2023, 10113601}, making it one of the fastest-growing consumer applications in history. While current precise user metrics for many LLMs are not consistently publicized, their societal impact has been increasing since then.

Following this surge in popularity, the LLM landscape has seen the emergence of numerous powerful proprietary models. Notable examples include Google's Gemini family of models, announced in December 2023 \cite{PichaiHassabis2023Gemini, Rane_Choudhary_Rane_2024}, and xAI's Grok, introduced in November 2023 \cite{xaiGrok2023}. These models, along with others like OpenAI's GPT-4 \cite{openai2023gpt4}, are considered at the time of this writing to be at the forefront of LLM performance, as evidenced by community benchmarks shown in the reference leaderboard Chatbot Arena \cite{chiang2024chatbotarenaopenplatform}. 
While exact parameter counts for these leading proprietary models are often not disclosed, they are understood to be exceptionally large, comparable in capabilities with models trained on trillions of parameters \cite{PichaiHassabisGemini2024}.

Concurrently, the open-source community has produced highly capable, albeit typically smaller, models. ``Commodity'' models, often in the range of 7 to 70 billion parameters, have demonstrated to be highly useful \cite{liu2023llm360fullytransparentopensource}. Prominent examples include the Mistral series, such as Mistral 7B \cite{jiang2023mistral7b}, and Meta's Llama family of models \cite{touvron2023llamaopenefficientfoundation}, which have fostered innovation and accessibility in the field.

The rise of LLMs has also begun to reshape the landscape of information retrieval.
Eventhough there is no official data, it is hinted by the concern within major search engine providers \cite{liu2024llm} that LLMs are capturing a notable share of tasks previously dominated by traditional search engines.
While not a complete replacement, LLMs offer a conversational paradigm for information synthesis and generation that complements, and in some cases competes with, established search methodologies.


\todo{Which countries use them the most according to ChatGPT? Which native languages do these countries have? How many speakers on English do they have? Are the countries that use it less partly because of their underrepresented language?}
Taking ChatGPT usage as a main example, due to it being the most popular LLM, a notable geographical concentrations in user activity can be seen. As of early 2025, web traffic data indicates that the United States accounts for the largest share of ChatGPT users (approximately 17\%), followed by India (around 8\%), and Brazil (approximately 5\%) \cite{SemrushChatGPTStatsApril2025}. Other countries such as Germany and France also feature among those with significant user bases \cite{SemrushChatGPTStatsApril2025, Views4YouChatGPTGlobal}.
An examination of the linguistic context within these leading user countries reveals important patterns. 
The United States primarily uses English, which is its de facto national language. 
India has English as an additional official language, widely employed in governance, higher education, and the technology sector, alongside Hindi and numerous other scheduled languages. 
In contrast, the primary language in Brazil is Portuguese, in Germany it is German, and in France it is French. 
English proficiency varies across these nations where English is not the primary language; for instance, Germany demonstrates ``High Proficiency'' in English, France ``Moderate Proficiency,'' while Brazil and India are categorized as having ``Low Proficiency'', according to the EF English Proficiency Index 2024 \cite{EFEPI2024}.

This distribution of usage, particularly the significant uptake in nations with varying levels of English proficiency and different primary languages, occurs in a context where many prominent LLMs are predominantly trained on English-language data \cite{zhang2023donttrustchatgptquestion, hao2024mindthegap}. This English-centricity in training data can lead to disparities in model performance across languages, often with suboptimal results for non-English, and especially low-resource, languages \cite{li2025fauxpolyglot}. The strong adoption in countries like India, despite overall lower English proficiency nationally, may reflect usage by specific English-speaking demographics or reliance on the LLMs' developing multilingual capabilities. Nevertheless, this situation highlights a potential linguistic divide, where users whose primary languages are underrepresented in LLM training data may not experience the full benefits of these technologies, or may encounter biases and reduced performance. Addressing this linguistic imbalance and improving the multilingual capabilities of LLMs are therefore critical areas of ongoing research and development.

\todo{introduce multilinguality, efforts in making LLMs multilingual, challenges, cross lingual transferability}

\todo{introduce high level how these are trainined}

\todo{introduce semantic web as an old concept. use right now? why do we need access to them? How could llms help us with different source languages}

\todo{introduce multilingual pretraining}
\todo{what is the aim of the thesis. what is some sneak peak, saying it is nuanced and not so straight forward as more languages more performance in every language}


